% ============================================================================
\chapter{Analysis and Properties}
\label{ch:analysis}
% ============================================================================

This chapter analyses the hybrid architecture from multiple
angles: accuracy, computational cost, energy efficiency,
interpretability, and theoretical properties.


% ----------------------------------------------------------------------------
\section{Accuracy Analysis}
\label{sec:accuracy-analysis}
% ----------------------------------------------------------------------------

\begin{table}[htbp]
\centering
\caption{Expected accuracy at each stage and comparison with
  baselines.}
\label{tab:accuracy-analysis}
\begin{tabular}{@{}lcl@{}}
  \toprule
  \textbf{Model} & \textbf{Test accuracy} & \textbf{Notes} \\
  \midrule
  TM alone (simple threshold) & 96.5\% & 500 clauses/class \\
  TM + LGN (continuous) & 97.5--98.0\% & Before discretisation \\
  TM + LGN + LUT (final) & 97.3--97.8\% & After quantisation \\
  \midrule
  2-layer FP network & 98.0\% & 784-300-10, ReLU \\
  Logistic regression & 92.0\% & Linear baseline \\
  Random forest & 97.0\% & Non-neural baseline \\
  LeNet-5 CNN & 99.0\% & Convolutional FP \\
  \bottomrule
\end{tabular}
\end{table}

The hybrid achieves accuracy comparable to a standard 2-layer FP
network (within 0.5\%) while using fundamentally different
operations. The remaining gap comes from two sources:

\begin{enumerate}
  \item \textbf{Booleanisation loss:} Thresholding pixels discards
    grayscale information. Thermometer encoding ($L = 4$) would
    recover $\sim$0.5\% accuracy at the cost of 4$\times$ larger input.

  \item \textbf{Discretisation loss:} The LGN's softmax relaxation
    is more expressive than the hardened circuit. Temperature
    annealing minimises but does not eliminate this gap.
\end{enumerate}


% ----------------------------------------------------------------------------
\section{Computational Cost Comparison}
\label{sec:cost-comparison}
% ----------------------------------------------------------------------------

\begin{table}[htbp]
\centering
\caption{Operation counts per image for different architectures.}
\label{tab:op-counts}
\begin{tabular}{@{}lrrr@{}}
  \toprule
  \textbf{Architecture} & \textbf{FP MACs} & \textbf{Bitwise ops}
    & \textbf{LUT reads} \\
  \midrule
  2-layer FP (784-300-10) & 238,200 & 0 & 0 \\
  Our hybrid & 0 & $\sim$252,000 & $\sim$16,010 \\
  ULEEN~\cite{susskind2024uleen} & 0 & 0 & $\sim$15,000 \\
  \bottomrule
\end{tabular}
\end{table}

The hybrid has a comparable number of operations to the FP baseline,
but each operation is $10$--$40\times$ cheaper in energy.


% ----------------------------------------------------------------------------
\section{Energy Efficiency}
\label{sec:energy}
% ----------------------------------------------------------------------------

Using the energy figures from \cref{tab:energy-cost}
(Horowitz~\cite{horowitz2014energy}, 45\,nm technology):

\paragraph{Floating-point baseline.}
$238{,}200$ FP MACs $\times$ $(3.7 + 0.9)$ pJ
$\approx 1{,}096{,}000$ pJ $\approx 1.1 \; \mu$J per image.

\paragraph{Hybrid architecture.}
\begin{itemize}
  \item Bitwise ops: $252{,}000 \times 0.1$ pJ $= 25{,}200$ pJ
  \item Small LUT reads (4-entry, register): $16{,}000 \times 0.1$
    pJ $\approx 1{,}600$ pJ
  \item Score LUT reads (L1 cache): $10 \times 5.0$ pJ $= 50$ pJ
  \item Popcount + adds: $5{,}000 \times 0.9$ pJ $= 4{,}500$ pJ
\end{itemize}
Total: $\sim$31,000 pJ $= 0.031 \; \mu$J per image.

\begin{keyinsight}
The hybrid architecture is approximately 35$\times$ more energy
efficient than the floating-point baseline for the same task. This
factor comes from two sources: cheaper operations (logic vs.\ FP
multiply) and smaller model (fits in cache vs.\ requires DRAM).
\end{keyinsight}

These are theoretical estimates based on operation energy costs~\cite{horowitz2014energy,wheeldon2020pervasive}.
Actual energy savings depend on the hardware platform, memory
hierarchy, and implementation efficiency. Real-world speedups of
$10$--$30\times$ are more realistic due to overhead.


% ----------------------------------------------------------------------------
\section{Memory and Cache Behaviour}
\label{sec:cache}
% ----------------------------------------------------------------------------

\begin{table}[htbp]
\centering
\caption{Model size breakdown and cache residency.}
\label{tab:model-size}
\begin{tabular}{@{}lrl@{}}
  \toprule
  \textbf{Component} & \textbf{Size} & \textbf{Fits in} \\
  \midrule
  Clause bitmasks & 980 kB & L2 cache \\
  Gate types & 8 kB & L1 cache \\
  Gate wiring & 26 kB & L1 cache \\
  Scoring LUTs & 4 kB & L1 cache \\
  \midrule
  \textbf{Total} & \textbf{1,018 kB} & L2 cache \\
  \bottomrule
\end{tabular}
\end{table}

The clause bitmasks are the dominant component (96\% of model size).
They fit in L2 cache (typically 256 kB--8 MB on modern CPUs) but not
in L1 (typically 32--64 kB). The LGN and LUT components easily fit
in L1.

\begin{engineernote}
For embedded deployment on microcontrollers with limited memory, the
clause bitmasks can be compressed. Sparse clauses (those with few
included literals) can be stored as lists of included literal indices
rather than full bitmasks, potentially reducing the TM component to
100--200 kB. This requires a different evaluation strategy (iterate
over included literals instead of bitwise AND) but saves memory.
\end{engineernote}


% ----------------------------------------------------------------------------
\section{Interpretability}
\label{sec:interpretability}
% ----------------------------------------------------------------------------

Each component of the hybrid architecture has different
interpretability properties:

\paragraph{Tsetlin Machine~\cite{granmo2018tsetlin} (high interpretability).}
Each clause is a conjunction of pixel conditions. A human can inspect
a clause and understand exactly what pattern it detects: ``pixel (14, 10)
must be bright AND pixel (5, 20) must be dark AND \ldots''. Clause
visualisation (\cref{fig:clause-vis}) makes this concrete.

\paragraph{Logic Gate Network~\cite{petersen2022difflogic} (medium interpretability).}
Each gate is a named Boolean function (AND, OR, XOR, etc.), and its
two inputs are identifiable. In principle, one could trace the logic
from output to input. In practice, with 16,000 gates in 4 layers,
this is difficult for humans but straightforward for automated
analysis tools.

\paragraph{LUT scoring (high interpretability).}
Each scoring table is a simple function mapping popcount to score.
It can be plotted and inspected directly. The shape reveals how
confidence relates to the number of firing gates.

\begin{keyinsight}
The hybrid architecture is substantially more interpretable than a
conventional neural network. A standard network's weights are opaque
floating-point numbers; the hybrid's parameters are readable Boolean
conditions, named logic gates, and plottable scoring curves. This is
valuable for debugging, auditing, and building trust in the model.
\end{keyinsight}


% ----------------------------------------------------------------------------
\section{Theoretical Properties}
\label{sec:theory}
% ----------------------------------------------------------------------------

\paragraph{Expressiveness.}
The LGN~\cite{petersen2024difflogic} with all 16 gate types is functionally complete
(\cref{sec:completeness}): it can represent any Boolean function of
its inputs given sufficient depth and width. The Tsetlin Machine~\cite{granmo2018tsetlin}
produces clauses that can approximate any Boolean function of the
pixels (since any function can be expressed as a disjunction of
conjunctions, i.e., a sum of clauses). Together, the architecture
is a universal Boolean function approximator.

\paragraph{Convergence.}
The Tsetlin Machine has formal convergence
guarantees~\cite{zhang2021convergence}: with sufficient states $N$,
each automaton converges to the optimal action. The LGN training
inherits the convergence properties of gradient descent on
differentiable functions.

\paragraph{Robustness.}
The Boolean nature of the architecture provides some inherent
robustness: small perturbations to pixel values (below the threshold)
do not change the booleanised input. However, adversarial examples
that flip pixels across the threshold can still fool the model.


% ----------------------------------------------------------------------------
\section{Limitations}
\label{sec:limitations}
% ----------------------------------------------------------------------------

\begin{enumerate}
  \item \textbf{Accuracy ceiling.} The $\sim$97.5\% accuracy on MNIST~\cite{lecun1998mnist} is
    below what CNNs achieve (99\%+). The gap comes from the fixed
    booleanisation and the lack of spatial structure (no convolution).

  \item \textbf{Scalability.} The Tsetlin Machine's clause count
    grows with task complexity. For larger images (e.g., CIFAR-10),
    many more clauses and literals are needed, and training time
    increases quadratically.

  \item \textbf{Two-phase training.} The sequential training
    (TM first, then LGN) means the TM cannot benefit from the LGN's
    feedback. Joint training would be ideal but is difficult because
    the TM and LGN use different learning paradigms.

  \item \textbf{Random connectivity.} The LGN's random wiring is
    simple but suboptimal. Learned or structured connectivity could
    improve accuracy but complicates the architecture.

  \item \textbf{No spatial awareness.} Unlike CNNs, the flat
    architecture treats every pixel independently. It cannot exploit
    the fact that nearby pixels are correlated. The Convolutional
    Tsetlin Machine~\cite{granmo2019convtm} addresses this but adds
    complexity.
\end{enumerate}


% ----------------------------------------------------------------------------
\section{Summary}
\label{sec:ch13-summary}
% ----------------------------------------------------------------------------

The hybrid architecture achieves its design goals:
\begin{itemize}
  \item \textbf{Accuracy:} 97--98\% on MNIST, competitive with
    FP baselines.
  \item \textbf{Efficiency:} 35$\times$ theoretical energy reduction
    vs.\ FP networks.
  \item \textbf{Simplicity:} Zero FP operations at inference. All
    operations are bitwise logic, integer math, and table lookups.
  \item \textbf{Interpretability:} Clauses, gates, and scoring
    tables are all human-inspectable.
  \item \textbf{Model size:} $\sim$1 MB, fits in L2 cache.
\end{itemize}

These properties come with trade-offs: lower accuracy than deep CNNs,
limited scalability to complex tasks, and sequential training that
prevents joint optimisation. The next chapter discusses how to
implement this architecture in Rust.
